\subsection{Basic properties}

\subsubsection{Definitions}\label{sssec:defs}

Let $S$ be a subset of places of ${\mathbb{Q}}$ and let
$||...||:G({\mathbb{A}}_S)\rightarrow {\mathbb{R}}_+$ be a ``height'' function
as in Section~\ref{sec:background}. Let $N>0$ and let
$\Gamma\subset \Pi(\overline{K})$
be a finite set of inequivalent genuine irreducible representations
of the n-fold cover of $K=K_{0,S}=\prod_{p\in S}K_{0,p}$.

We call a genuine factorizable function
$f=\prod_{p\in S}f_p:\overline{G({\mathbb{A}}_S)}\rightarrow {\mathbb{C}}$
{\bf smooth} if

(a) $f_\infty\in C^\infty(\overline{G({\mathbb{R}})})_{gen}$
if $p=\infty \in S$, and

(b) $f_p\in C_c^\infty(\overline{G({\mathbb{Q}}_p)})_{gen}$
if $p\in S-\{\infty\}$.

\noindent
In general, a genuine
function $f:\overline{G({\mathbb{A}}_S)}\rightarrow {\mathbb{C}}$
is {\bf smooth} if it is a finite linear combination of factorizable smooth
functions. Define
${\mathcal{H}}_N(\overline{G({\mathbb{A}}_S)})_\Gamma$
to be the vector space of
all genuine smooth complex valued functions
$f:\overline{G({\mathbb{A}}_S)}\rightarrow {\mathbb{C}}$
such that
\begin{equation}\label{eq:4.11}
\begin{array}{c}
{\rm supp}(f)\subset \overline{G({\mathbb{A}}_S,N)}
=\{ g\in \overline{G({\mathbb{A}}_S)}\ |\ ||g||\leq e^N\}\\
{\rm The \ spans\ of\ }
k\longmapsto f(kg){\rm \ generates\ a \ repn\ in\ }\Gamma,\\
{\rm The \ spans\ of\ }
k\longmapsto f(gk){\rm \ generates\ a \ repn\ in\ }\Gamma .
\end{array}
\end{equation}
The semi-norms
\[
||f||_D=\sup_{x\in \overline{G({\mathbb{A}}_S)}} |Df(x)|,
\]
$D$ a differential operator on $\overline{G({\mathbb{R}})}$,
induces a topology on
${\mathcal{H}}_N(\overline{G({\mathbb{A}}_S)})_\Gamma$
making it into a topological
vector space. By \cite{Ru}, ch 1,
we may complete this space using the above
family of seminorms and obtain a Frech\'et space.

Define
\[
{\mathcal{H}}(\overline{G({\mathbb{A}}_S)})=\lim_{\Gamma}
{\mathcal{H}}(\overline{G({\mathbb{A}}_S)})_\Gamma
\]
where
\[
{\mathcal{H}}(\overline{G({\mathbb{A}}_S)})_\Gamma =\lim_{N}
{\mathcal{H}}_N(\overline{G({\mathbb{A}}_S)})_\Gamma .
\]
Therefore, each function in
${\mathcal{H}}(\overline{G({\mathbb{A}}_S)})$
is both right and left $\overline{K}$-finite.

Similarly, if $A$ denotes the usual
diagonal torus of $G=SL(2)$, then we define
${\mathcal{H}}(\overline{A({\mathbb{A}}_S)})$.

From \cite[section 11]{A11}, another space of functions of
``almost compact'' support,
${\mathcal{H}}_{ac}(\overline{G({\mathbb{A}}_S)})$, is introduced,
as well as its subspace of ``moderate'' functions,
which we denote by
${\mathcal{H}}_{mod}(\overline{G({\mathbb{A}}_S)})$.
However, in the case $G=SL(2)$ it turns out that
\[
\begin{array}{c}
{\mathcal{H}}_{mod}(\overline{G({\mathbb{A}}_S)})=
{\mathcal{H}}_{ac}(\overline{G({\mathbb{A}}_S)})=
{\mathcal{H}}(\overline{G({\mathbb{A}}_S)}),\\  
{\mathcal{H}}_{mod}(\overline{A({\mathbb{A}}_S)})=
{\mathcal{H}}_{ac}(\overline{A({\mathbb{A}}_S)})=
{\mathcal{H}}(\overline{A({\mathbb{A}}_S)}).
\end{array}
\]

We shall only sketch the definition of the Schwartz
space ${\mathcal{C}}(\overline{G({\mathbb{A}}_S)})$: this is the space spanned by
factorizable functions $f$ each of whose local components
belong to the local Schwartz space, defined for example
in \cite{J1}.
% Claude (2026): the 1997 book cites its Chapter 6 here.
This space will play a bigger role in Section~\ref{sec:invariant} below.

\subsection{An equivalence relation on $G({\mathbb{Q}})$}

Let $\gamma ,\gamma ^{\prime }\in G({\mathbb{Q}})$ and let
$\gamma =\gamma_u\gamma _s$ denote its Jordan decomposition.
We say that $\gamma ,\gamma^{\prime }$ are \textbf{equivalent},
written $\gamma \sim \gamma ^{\prime }$,
if $\gamma _s$ is $G({\mathbb{Q}})$-conjugate to
$\gamma _s^{\prime }$. Clearly
if $\gamma ,\gamma ^{\prime }$ are conjugate and
semi-simple then they are
equivalent. In fact, by the uniqueness of the Jordan decomposition, it can
be shown that if $\gamma ,\gamma ^{\prime }$ are conjugate then they are
equivalent. Define $\mathcal{C}$ to be the set of all equivalence classes.
This definition holds in general.

At least for $G=SL(n)$, the classical analog of this definition is the
following: Let $\gamma ,\gamma ^{\prime }\in G({\mathbb{Z}})$. We say that
$\gamma ,\gamma ^{\prime }$ are \textbf{equivalent}, written $\gamma \sim
\gamma ^{\prime }$, if $\gamma _s$ is $G(\mathbb{C})$-conjugate to $\gamma
_s^{\prime } $. In other words, $\gamma ,\gamma ^{\prime }$ are equivalent
if they have the same set of eigenvalues.

We call $(\gamma ,\zeta)\in \overline{G({\mathbb{A}})}$
semi-simple if $\gamma\in G({\mathbb{A}})$ is semi-simple.
Let ${\mathcal{O}}\in \mathcal{C}$ denote a
class consisting entirely of semisimple
elements (we call such a class \textbf{semi-simple}).
If $\gamma \in {\mathcal{O}}$
then the centralizer $G_\gamma $ is reductive and anisotropic modulo its
center:

\begin{equation}\label{eq:4.12}
\begin{array}{c}
G_\gamma =\left\{ 
\begin{array}{c}
G,\;\;\;\;\;\;\;\;\;\;\;\;\;\;\;\;\;\;\;\;\;\gamma =\pm 1, \\ 
A^x,\;\;\;\gamma \neq \pm 1\;and\;\gamma \sim a\in A, \\ 
T,\;\;\;\;\;\;\;\;\;\;\;\;\;\;\;\;\;\;\mathrm{otherwise},
\end{array}
\right.
\end{array}
\end{equation}
where $A^x=x^{-1}Ax$ and $x^{-1}\gamma x\in A$.
In particular, all such classes are unramified in the sense of
\cite{A2}, p. 949.

Let ${\mathcal{O}}\in \mathcal{C}$ denote an unramified
semi-simple class, let $P=MN\in\{B,G\}$, and let
$\mu \in {\mathcal{O}}\cap M({\mathbb{Q}})$.
Then

\begin{equation}\label{eq:4.13}
\begin{array}{c}
J_{B,{\mathcal{O}}}(\delta x,\delta x)
=\sum_{\eta \in M_\mu ( {\mathbb{Q}})\backslash M(\mathbb{Q})}
\sum_{\varsigma \in N({\mathbb{Q}})}
f(x^{-1}\delta ^{-1}\varsigma ^{-1}\eta^{-1}\mu \eta \varsigma \delta x) \\ 
=\sum_{\eta \in M_\mu ({\mathbb{Q}})\backslash P({\mathbb{Q}})}
f(x^{-1}\delta^{-1}\eta ^{-1}\mu \eta \delta x).
\end{array}
\end{equation}

If $P=B$ then $M=A$ and $M_\mu =A$ for
$\mu \in {\mathcal{O}}$ as above. If $P=G$ then
$M=G$ and $M_\mu =G_\mu $ is one of the above
centralizers for $\mu \in {\mathcal{O}}$ as
above. As we saw above, if $P=B$ then

\begin{equation}\label{eq:4.14}
\begin{array}{c}
J_{B,{\mathcal{O}}}(\delta x,\delta x)
=\sum_{\varsigma \in N( {\mathbb{Q}})}
f(x^{-1}\delta^{-1}\varsigma ^{-1}\mu \varsigma \delta x) \\ 
=\sum_{\eta \in A({\mathbb{Q}})\backslash B({\mathbb{Q}})}
f(x^{-1}\delta ^{-1}\eta^{-1}\mu \eta \delta x).
\end{array}
\end{equation}
We have then by an equation above

\begin{equation}\label{eq:4.15}
\begin{array}{c}
j_{\mathcal{O}}^T(x,f)
=\sum_{P\in \{B,G\}}(-1)^{\dim A_P}
\sum_{\delta \in P( {\mathbb{Q}})\backslash G({\mathbb{Q}})}
J_{P,{\mathcal{O}}}(\delta x,\delta x) \tau _P^{\wedge}(H_B(\delta x)-T) \\ 
=\sum_{P\in \{B,G\}}(-1)^{\dim A_P}
\sum_{\eta \in M_\mu ( {\mathbb{Q}})\backslash G({\mathbb{Q}})}
f(x^{-1}\eta ^{-1}\mu \eta x)\tau _P^{\wedge}(H_B(\eta x)-T) \\ 
=\sum_{\eta \in G_\mu ({\mathbb{Q}})\backslash G({\mathbb{Q}} )}
f(x^{-1}\eta ^{-1}\mu\eta x)
\sum_{P\in \{B,G\}}(-1)^{\dim A_P} \tau_P^{\wedge }(H_B(\eta x)-T).
\end{array}
\end{equation}
Also

\begin{equation}\label{eq:4.16}
\begin{array}{c}
J_{\mathcal{O}}^T(f)
=\int_{G({\mathbb{Q}})\backslash
\overline{G({\mathbb{A}})}}
j_{\mathcal{O}}^T(x)dx \\ 
=\int_{G({\mathbb{Q}})\backslash \overline{G({\mathbb{A}})}}
\sum_{\delta \in G_\mu ({\mathbb{Q}})\backslash G({\mathbb{Q}})}
f(x^{-1}\delta ^{-1}\mu \delta x)
\sum_{P\in\{B,G\}}(-1)^{\dim A_P}\tau _P^{\wedge }(H_B(\delta x)-T)dx \\ 
=\int_{G_\mu ({\mathbb{Q}})\backslash \overline{G({\mathbb{A}})}}
f(x^{-1}\mu x)\sum_{P\in\{B,G\}}(-1)^{\dim A_P}
\tau _P^{\wedge }(H_B(x)-T)dx \\ 
=vol(G_\mu ({\mathbb{Q}})\backslash \overline{G({\mathbb{A}})}_\mu)
\int_{\overline{G({\mathbb{A}})}_\mu\backslash \overline{G({\mathbb{A}})}}
f(x^{-1}\mu x)\sum_{P\in \{B,G\}}(-1)^{\dim A_P}
\tau_P^{\wedge }(H_B(x)-T)dx \\ 
=vol(A_{\mathcal{O}}({\mathbb{R}})^oG_\mu ({\mathbb{Q}})\backslash
\overline{G ({\mathbb{A}})}_\mu)
\int_{\overline{G ({\mathbb{A}})_\mu}\backslash \overline{G({\mathbb{A}})}}
f(x^{-1}\mu x)v(x,T)dx
\end{array}
\end{equation}
(see also \cite[(5.5)]{Sho}), where 
\begin{equation}\label{eq:4.17}
v(x,T)=\int_{A_{\mathcal{O}}({\mathbb{R}})^o}
\sum_{P\in \{B,G\}}(-1)^{\dim A_P}\tau_P^{\wedge }(H_B(ax)-T)da.
\end{equation}

If $A_{\mathcal{O}}=A$, that is to say if there is an
$a\in A({\mathbb{Q}})-\{\pm I_2\}$ in the class ${\mathcal{O}}$
then $v(x,T)$ is the ``volume'' of the convex hull of 
\[
\{T-H(x),-(T-H(wx))\}. 
\]
Thus 
\begin{equation}\label{eq:4.18}
v(x,T)=2T-\log \delta(x)-\log \delta(wx).
\end{equation}

This proves the following

\begin{theorem}\label{thm:4.10}
Let ${\mathcal{O}}\in \mathcal{C}$ and suppose that
${\mathcal{O}}\cap (A({\mathbb{Q}})-\{\pm I_2\})\not=\emptyset $.
Then 
\[
J_{\mathcal{O}}^T(f)
=vol(A_{\mathcal{O}}({\mathbb{R}})^oG_\mu ({\mathbb{Q}})\backslash G_\mu ({\mathbb{A}}))
\int_{G_\mu ({\mathbb{A}})\backslash G({\mathbb{A}})}f(x^{-1}\mu x)v(x,T)dx. 
\]
\end{theorem}

These are called \textbf{weighted orbital integrals}.
The case where $1\in {\mathcal{O}}$ will be discussed below.

\subsection{The definition of $J_M(u,f)$}\label{ssec:JMu}

Let $f\in C_c^\infty (\overline{G({\mathbb {A}}_S)})_{gen}$ and let
$\overline{\gamma}=(\gamma,1) \in \overline{M({\mathbb {A}}_S)}$
have Jordan decomposition
$\gamma =\sigma u$. Suppose that
$\overline{M({\mathbb{A}}_S)}_{\overline{\gamma}}
=\overline{G({\mathbb{A}}_S)}_{\overline{\gamma}} $.
Assume that $S$ is a ``sufficiently large'' finite set which contains
the archimedian places. In this case, define 
\begin{equation}\label{eq:4.19}
J_M(\overline{\gamma} ,f)=|D(\gamma )|^{1/2}
\int_{\overline{M({\mathbb {A}}_S)}
\backslash \overline{G({\mathbb {A}}_S)}}
\int_{{\mathcal O}_{\overline{\gamma}} (\overline{M({\mathbb {A}}_S)})}
f(x^{-1}\mu x)v_M(x)d\mu dx,
\end{equation}
where ${\mathcal O}_{\overline{\gamma}} (\overline{M({\mathbb {A}}_S)})
=\{m^{-1}\overline{\gamma} m\;m\in \overline{M({\mathbb {A}}_S)})$.
Here $v_M$ is the ``volume'' of the convex hull of the points
$\{H_B(x),H_{B^{-}}(x)\}$ in ${\mathfrak a}$. If $x=nk$ then
$H_B(x)=0$, so $v_M(nk)=|H_{B^{-}}(x)|$.
In any case, this integral can be written as 
\[
J_M(\overline{\gamma} ,f)=|D(\gamma )|^{1/2}
\int_{\overline{G({\mathbb {A}}_S)}_{\overline{\gamma}}
\backslash \overline{G({\mathbb{A}}_S)}}
f(x^{-1}\overline{\gamma} x)v_M(x)dx. 
\]
This is a function on $\overline{M({\mathbb {A}}_S)}$.

If $\gamma \in \overline{M({\mathbb {A}}_S)}$
is {\em arbitrary} and $a\in \overline{A_M({\mathbb {A}}_S)}$ is
chosen so that
$M({\mathbb {A}}_S)_{a\gamma }=G({\mathbb {A}}_S)_{a\gamma }$
then $J_M(a\gamma ,f)$ is well-defined. We may consider
$a\rightarrow 1$.

As in \cite[\S 3]{A7}, let 
\[
r_\beta (\lambda ,u,a)
=|a^\beta -a^{-\beta }|^{\rho (\beta ,u)\lambda (\beta^{\vee })},
\;\;\;\;\lambda \in {\mathfrak {a}}_{M,{\mathbb {C}}}^{*}, 
\]
and let 
\[
r_P(\lambda ,u,a)
=\prod_{v\in S}\prod_\beta
r_\beta (\frac \lambda 2,u_v,a_v), 
\]
where $\beta $ runs over the reduced roots of
$(P_{\sigma _v},A_{M_{\sigma_v}})$. By definition, 
\[
r_M(u ,a)=\lim_{\lambda \rightarrow 0}
\sum_{P\in {\mathcal P}(M)}r_P(\lambda ,u,a)\theta _P(\lambda )^{-1}, 
\]
and 
\[
v_M(x)=\lim_{\lambda \rightarrow 0}
\sum_{P\in {\mathcal P}(M)}e^{-\lambda (H_P(x))}
\theta _P(\lambda )^{-1}. 
\]
Following \cite[(8.1)]{A7},
we define $J_M(\gamma ,f)$ to be the limit at $a=1$
of 
\begin{equation}\label{eq:4.20}
\begin{array}{c}
|D(a\gamma )|^{1/2}
\int_{\overline{G({\mathbb {A}}_S)}_\sigma \backslash
\overline{G({\mathbb {A}}_S)}}
\int_{\overline{M({\mathbb {A}}_S)}_\gamma
\backslash \overline{G({\mathbb {A}}_S)}_\sigma }
f(y^{-1}\sigma x^{-1}a u xy)\times \\
\times
\left( \sum_{L\in {\mathcal{L}}(M)}
r_M^L(\gamma ,a)v_L(xy)\right) dxdy,
\end{array}
\end{equation}
where $\gamma\in \overline{M({\mathbb A}_S)}$ is as above.
\begin{example}\label{ex:4.11}
Let $G=SL(2),\;M=A,\;a=\left( 
\begin{array}{cc}
t & 0 \\ 
0 & t^{-1}
\end{array}
\right) $, where $t\neq \pm 1$. In this case, 
\[
r_A(\gamma ,a)=1,
\]
and 
\[
v_A(x)=\sum_{p\in S}\log (\sqrt{|1+n_p^2|_p}),
\]
where $x=\prod_{p\in S}a_p\cdot \left[ 
\begin{array}{cc}
1 & n_p \\ 
0 & 1
\end{array}
\right] \cdot k_p,$ for
$a_p\in \overline{A({\mathbb {Q}}_p)}$, $k_p\in \overline{K_{0,p}}$
(\cite{A6}, Intro., \cite{Sho}, chap. 6). Therefore,
$J_A(z,f)$, for $z\in Z(G({\mathbb Q}))$,
is the limit as $a\rightarrow 1$ of 
\begin{equation}\label{eq:4.21}
\ |D(a)|^{1/2}
\int_{\overline{A({\mathbb {A}}_S)}\backslash \overline{G({\mathbb {A}}_S)}}
f(y^{-1} az y)v_A(y)dy 
\end{equation}
\cite[section 6.1]{Sho}, p. 64.
\end{example}

\subsection{A more explicit form of the geometric terms}\label{ssec:explicit}

\begin{definition}\label{defn:4.12}
Let 
\[
K_{P,unip}(x,y)
=\sum_{\gamma \in {\mathcal U}_{M_P}{\mathbb {(Q)}}}
\int_{N_P({\mathbb {A})}}f(x^{-1}\gamma ny)dn,
\]
and for $T_0=0$, let 
\begin{eqnarray*}
J_{unip}(f)
&=&J_{unip}^{T_0}(f) \\
&=&\int_{G({\mathbb{Q}})\backslash \overline{G({\mathbb{A}})}}
\sum_P(-1)^{\dim(A_P/A_G)}
\sum_{\delta \in P({\mathbb {Q}})\backslash G({\mathbb{Q}})}
K_{P,unip}(\delta x,\delta x)
\tau _P^{\wedge }(H(\delta x)-T_0)dx.
\end{eqnarray*}
\end{definition}

\begin{example}\label{ex:4.13}
Let $G=SL(2),\;M=A$. Then 
\[
K_{P,unip}(x,y)
=\sum_{\gamma \in {\mathcal U}_{M_P}({\mathbb {Q}})}
\int_{N_P({\mathbb {A})}}f(x^{-1}\gamma ny)dn
\]
and 
\begin{eqnarray*}
J_{unip}(f)
=\int_{G({\mathbb {Q}})\backslash \overline{G({\mathbb{A}})}}
\sum_P(-1)^{\dim(A_P/A_G)}
\sum_{\delta \in P({\mathbb {Q}})\backslash G({\mathbb{Q}})}
K_{P,unip}(\delta x,\delta x)
\tau _P^{\wedge }(H(\delta x)-T_0)dx \\
=\int_{G({\mathbb{Q}})\backslash \overline{G({\mathbb{A}})}}
(K_{G,unip}(x,x)-
\sum_{\delta \in B({\mathbb {Q}})\backslash G({\mathbb{Q}})}
K_{B,unip}(\delta x,\delta x)
\tau _B^{\wedge }(H(\delta x)-T_0))dx
\end{eqnarray*}
\end{example}

\begin{definition}\label{defn:4.14}
Let $J_M(u,f)$ denote the weighted unipotent orbital integrals defined in
the previous subsection,
$u\in M({\mathbb A}_S)$ unipotent.
We define $a^M(S,u)$ by the property that 
\[
J_{unip}(f)
=\sum_P
\sum_{u\in {\mathcal U}_{M_P}({\mathbb {Q}})}
|W_0^M|\cdot |W_0^G|^{-1}a^M(S,u)J_M(u,f).
\]
(The verification that this definition makes sense can be found
for example in \cite[section 6.2, pp.~66--67]{Sho}.) For
$\gamma\in \overline{M({\mathbb A}_S)}$, we define $a^M(S,\gamma )$ by 
\[
a^M(S,\gamma )
=\varepsilon ^M(\sigma )|\iota ^M(\sigma)|^{-1}
\sum_{\{u\;|\;\sigma u\sim \gamma \}}a^{M_\sigma }(S,u),
\]
where 
\[
\varepsilon ^M(\sigma )=\left\{ 
\begin{array}{c}
1,\;\;{\rm if}\ \sigma \;{\rm is\ }k-{\rm elliptic\;on}\;M \\ 
0,\;\;\;{\rm otherwise}.\;\;\;\;\;\;\;\;
\end{array}
\right. 
\]
Furthermore, in the case of $G=SL(2)$, if $\gamma =\sigma u$
is the Jordan decomposition then
$a^M(S,\gamma)=\epsilon^M(\sigma)a^M(S,u)$.
\end{definition}

The relationship between the $J_{{\mathcal O}}(f)$ and the
$J_M(f)$ is given by the following:

\begin{proposition}\label{prop:4.15}
(\cite[Theorem 8.1]{A12}) If $S$ is a sufficiently large finite set of
valuations and
$f\in C_c^\infty (\overline{G({\mathbb {A}}_S)})_{gen}$ then 
\[
J_{{\mathcal O}}(f)=\sum_{M\in \{A,G\}}
|W_0^M|\cdot |W_0^G|^{-1}
\sum_{\gamma \in M({\mathbb {Q}})\cap {\mathcal{O}}}
a^M(S,\gamma )J_M(\gamma ,f) 
\]
where the constants $a^M(S,\gamma )$ defined above
satisfy, if $\gamma $ is semisimple, the property 
\[
a^M(S,\gamma )=\left\{ 
\begin{array}{c}
vol(M({\mathbb {Q}})_\gamma \backslash \overline{M({\mathbb {A}}_S)}_\gamma ^1),
\;\;\;\;{\rm if\;}\gamma \;{\rm is}\;
{\mathbb {Q}}-{\rm elliptic\;in}\;M({\mathbb {Q}}), \\ 
0,\;\;\;\;\;\;\;\;\;\;\;\;\;\;\;\;\;\;{\rm otherwise},\;\;\;\;\;
\end{array}
\right. 
\]
\end{proposition}

\begin{example}\label{ex:4.16} For $G=SL(2)$, $1\in {\mathcal O}$, we have
\[
J_{\mathcal O}(f)=\sum_{\gamma\in {\mathcal{O}}}
a^G(S,\gamma)J_G(\gamma,f)
+{1\over 2}a^A(S,1)J_A(1,f)
\]
where $J_A(1,f)$ is in Example~\ref{ex:4.13} above and $J_G(\gamma,f)$
is defined in \eqref{eq:4.19} above.
A similar result holds for the ${\mathcal{O}}$ with $-1\in {\mathcal{O}}$.
\end{example}


